\section{Conclusions}
This paper followed one argument through six steps on the public data of Nishimura et al. (1)~The two-step semi-mechanistic approach is transparent and reproducible from its published data to the last digit for the data checks, the fingerprint step and the simplified dynamic fits; but several printed full-model fits are not at their optimum, the full model is structurally non-identifiable, and the Tao case study cannot be reproduced, so in-sample fit statistics are not a sufficient basis for model selection. (2)~The misfit the method leaves has a mechanistic signature: pseudo-processivity in the mass-transfer-limited Bornadel data and, specific to a single profile, within-batch deactivation at 100\,°C. (3)~PPK-FIP keeps the Ping Pong core and the idea that selectivities are best learned against the degree of reaction, but uses the fingerprint estimate as prior information, adds the two mechanisms as testable options, and is fitted globally with an exact, fast time reparameterisation; it fits all six case studies better and predicts held-out observations with \LotoAllRedOrig\,\% lower error than the original method and \LotoAllRedFull\,\% lower than its full variant. (4)~Pseudo-processivity and the calibrated prior carry this gain; deactivation is a profile-specific option that the prior activates only where a profile slows down late, and the full lumped-constant set can replace model averaging. (5)~With known truth, the two-step method overestimates the selectivity $S_M$ by \SimSMBiasOrigPct\,\% when pseudo-processivity is present, whereas PPK-FIP is nearly unbiased and reduces to the original model when processivity is absent; the processive path transfers across the Tao water activities. (6)~Black-box learners interpolate dense profiles but fail in extrapolation; combined with a Gaussian-process discrepancy (PPK-FIP-GP), the mechanistic model halves the cross-validated error of the original method and outperforms every black-box model while extrapolating like the mechanism.

For the design of biolubricant processes from biomass-derived fatty acids and TMP this means that apparent selectivities obtained in mass-transfer-limited reactors should not be transferred to other reactor configurations without a processivity correction, that the fingerprint idea is most valuable as prior information rather than as a constraint, and that kinetic parameters should be reported with out-of-sample validation rather than in-sample fit statistics alone.
